\section*{Bab \ref{ch:b2:lab-techniques-2} --- Teknik Laboratorium II: Sintesis dan Analisis dalam Praktik}

\begin{solution}{exo:b2:lab-techniques-2:1}
Es dan air; es dan garam; es kering dan aseton (seperti untuk \omterm{def:b2:enolates-aldol:enolate}{enolat} pada \cref{ch:b2:enolates-aldol}).
\end{solution}

\begin{solution}{exo:b2:lab-techniques-2:2}
Lapisan bawah. Tambahkan beberapa tetes air: tetes-tetes itu bergabung dengan lapisan berair, di sini lapisan atas.
\end{solution}

\begin{solution}{exo:b2:lab-techniques-2:3}
Porsi-porsi pertama menggumpal karena menyerap air; ketika padatan yang baru
ditambahkan tetap lepas-lepas dan berputar dengan bebas, larutannya sudah
kering.
\end{solution}

\begin{solution}{exo:b2:lab-techniques-2:4}
$R_f$ terlalu tinggi untuk kolom (yang diinginkan sekitar 0.3) dan kedua
bercak berdekatan: pakailah eluen yang kurang polar, lebih banyak heksana
(misalnya 9\,:\,1), yang menurunkan kedua $R_f$ dan memperlebar jarak
keduanya dalam volume kolom.
\end{solution}

\begin{solution}{exo:b2:lab-techniques-2:5}
\begin{itemize}
  \item Natrium hidrogenkarbonat (pH 8.3 $>$ 4.2 + 2): benzoat masuk ke air;
        asamkan, ekstraksi asam benzoat.
  \item Natrium hidroksida (pH $>$ 12 $>$ 9.99 + 2): fenoksida masuk ke air;
        asamkan, ekstraksi fenol.
  \item Asam klorida (pH sekitar 1 $<$ 4.6 $-$ 2): anilinium masuk ke air;
        basakan, ekstraksi anilina.
  \item Naftalena tetap berada dalam eter: keringkan, uapkan.
\end{itemize}
Pencucian dengan hidrogenkarbonat harus lebih dahulu: hidroksida akan mengekstraksi asam dan fenol bersama-sama.
% ledger: ph:NaHCO3, pka:benzoic, pka:phenol, pka:anilinium
\end{solution}

\begin{solution}{exo:b2:lab-techniques-2:6}
$1/T = 1/T_b - (R/\Delta_{\mathrm{vap}}H)\ln(p/p^\circ)$: bromobenzena sekitar \qty{49}{\degreeCelsius},
benzaldehida sekitar \qty{68}{\degreeCelsius} (gambar pada \cref{prop:b2:lab-techniques-2:reduced-pressure}).
\end{solution}

\begin{solution}{exo:b2:lab-techniques-2:7}
$0.050 \times 200\,000/400 = \qty{25}{K}$.
\end{solution}

\begin{solution}{exo:b2:lab-techniques-2:8}
Setiap senyawa mempunyai tanggapannya sendiri: persen luas adalah persen
massa hanya jika faktor-faktor responsnya sama. NMR: $0.12/3 = 0.040$~mol
pengotor per mol produk, yaitu $0.040 \times 100 = \qty{4}{g}$
per \qty{200}{g} produk; $w = 200/204 \approx 98.0~\%$.
\end{solution}

\begin{solution}{exo:b2:lab-techniques-2:9}
Massa teoretis $0.0300 \times 164 = \qty{4.92}{g}$. Tanpa koreksi $4.10/4.92 \approx 83.3~\%$; dengan koreksi
$0.950 \times 4.10/4.92 \approx 79.2~\%$.
\end{solution}

\begin{solution}{exo:b2:lab-techniques-2:10}
Benzena berarti pereaksi itu terbentuk lalu terprotonasi: air di dalam alat
gelas, pelarut, atau aldehida; udara (kelembapan) yang masuk melalui
kebocoran; aldehida yang sebagian teroksidasi menjadi asam benzoat, yang
hidrogen asamnya merusak pereaksi. Penanggulangannya: alat gelas yang
dikeringkan dalam oven, pelarut kering, aldehida yang baru didistilasi,
tekanan nitrogen berlebih yang diperiksa pada penggelembung. Jika hampir
tidak ada pereaksi sama sekali, itu menunjukkan kegagalan inisiasi
(permukaan magnesium yang tertutup oksida), yang dihindari dengan serutan
yang segar atau diaktifkan.
\end{solution}

\begin{solution}{exo:b2:lab-techniques-2:11}
\qty{270}{kJ} pada \qty{150}{W}: sekurang-kurangnya $270\,000/150 = \qty{1800}{s}$,
setengah jam, dan hanya jika halida itu bereaksi secepat penambahannya.
Jika lebih cepat, pelepasan kalor melampaui pendinginan: suhunya naik, atau
pereaksi terakumulasi dan kemudian dapat bereaksi sekaligus
(\cref{prop:b2:lab-techniques-2:accumulation}).
\end{solution}

\begin{solution}{exo:b2:lab-techniques-2:12}
Rekristalisasi: $13.0 \times 0.80 = \qty{10.4}{g}$ pada 99~\%. Kolom: $13.0 \times 0.95 =
\qty{12.4}{g}$ pada 99.5~\%, tetapi dengan \qty{1.5}{L} pelarut yang harus
dibeli, diuapkan, dan dibuang. Kolom memberikan sekitar \qty{2}{g} produk
lebih banyak yang sedikit lebih murni; rekristalisasi lebih sederhana,
lebih murah, dan lebih sedikit limbahnya, dan lebih disukai jika
kemurniannya sudah cukup.
\end{solution}

\begin{solution}{pb:b2:lab-techniques-2:1}
\textbf{1.} Etoksietana: sangat mudah terbakar, berbahaya, membentuk
peroksida. Bromobenzena: mudah terbakar, beracun, beracun bagi kehidupan
akuatik. Magnesium: padatan mudah terbakar, melepaskan hidrogen dengan air.
\textbf{2.} \ce{C6H5MgBr + H2O -> C6H6 + Mg(OH)Br}: setiap tetes air merusak
pereaksi; pengeringan dalam oven menyingkirkan lapisan tipis air pada kaca.
\textbf{3.} Untuk mencegah masuknya kelembapan dan oksigen, yang juga merusak pereaksi.
\textbf{4.} Aliran nitrogen; penggelembung mencegah udara masuk melalui saluran keluar.
\textbf{5.} Etoksietana tidak bereaksi dengan pereaksi, pasangan elektron
bebas oksigennya menstabilkan magnesium, dan titik didihnya yang rendah
(\qty{307.7}{K}) memungkinkan refluks membatasi suhu.
\textbf{6.} $15.7/157.01 = \qty{0.1000}{mol}$ bromobenzena; $2.67/24.305 = \qty{0.1099}{mol}$ magnesium,
berlebih 10~\%.
\textbf{7.} $0.1000 \times 270 = \qty{27.0}{kJ}$.
\textbf{8.} Agar kalor dilepaskan dengan laju yang dapat disingkirkan oleh
refluks dan penangas.
\textbf{9.} $0.0200 \times 270 = \qty{5.40}{kJ}$.
\textbf{10.} $5400/170 \approx \qty{32}{K}$.
\textbf{11.} Sekitar \qty{52}{\degreeCelsius}, jauh di atas titik didih
etoksietana (\qty{34.6}{\degreeCelsius}): pendidihan hebat yang dapat
membanjiri kondensor dan melemparkan uap yang mudah terbakar keluar dari
labu.
\textbf{12.} Tambahkan sekitar sepersepuluh bromobenzena, tunggu sampai
reaksi dimulai (kekeruhan, refluks yang lembut), lalu tambahkan sisanya
dengan laju yang mempertahankan refluks lembut, dengan penangas es yang
siap.
\textbf{13.} \ce{C6H5MgBr + C6H5CHO} menghasilkan alkoksida \ce{(C6H5)2CHOMgBr}; benzaldehida,
$9.55/106.12 = \qty{0.0900}{mol}$, menjadi pembatas.
\textbf{14.} Alkohol itu dua kali benzilik: asam kuat akan mengubahnya
menjadi kation yang terstabilkan, yang mengarah ke eter atau produk
substitusi; amonium klorida cukup asam untuk memprotonasi alkoksida dan
melarutkan garam-garam magnesium.
\textbf{15.} Lapisan atas: etoksietana kurang rapat daripada air.
\textbf{16.} Pencucian itu menyingkirkan sebagian besar air yang terlarut dalam eter.
\textbf{17.} Magnesium sulfat sampai lepas-lepas, penyaringan, lalu penguap
putar pada tekanan rendah.
\textbf{18.} Fenilmagnesium bromida berkopling dengan bromobenzena yang belum
bereaksi; konsentrasi bromobenzena lokal yang tinggi dan suhu tinggi
mendukungnya, alasan lain untuk penambahan yang lambat.
\textbf{19.} Bifenil ($R_f$ 0.85), yang kurang polar.
\textbf{20.} Sekitar $1/0.85 \approx 1.2$ volume kolom untuk bifenil, $1/0.25 = 4$ untuk alkohol.
\textbf{21.} Alkohol menyumbang 10 proton \omterm{def:b2:huckel:aromatic}{aromatik} per CH; kelebihan
$10.90 - 10.00 = 0.90$ berasal dari bifenil, 10 proton per molekul: 0.090
mol bifenil per mol alkohol.
\textbf{22.} $w = 184.24/(184.24 + 0.090 \times 154.21) \approx 0.930$.
\textbf{23.} $12.50 \times 0.930 \approx \qty{11.62}{g}$.
\textbf{24.} $0.0900 \times 184.24 \approx \qty{16.58}{g}$.
\textbf{25.} $12.50/16.58 \approx 75.4~\%$.
\textbf{26.} $11.62/16.58 \approx \boldsymbol{70~\%}$ (70.1~\%), dikoreksi untuk kemurnian.
\end{solution}
